\chapter{根与翼} % Introduction chapter suppressed from the table of contents

中国社会一直非常强调家庭价值观，希望实现家族的持续传承，家族有族谱，代代相传的关系对每个家庭成员的成长产生深远影响。我们每个人都只是人类进化过程中的短暂过渡。父母普遍希望把最好的东西传承给下一代。然而我们需要问自己，什么东西能够持久，仅仅是财富和金钱吗？\\
继承大量财富，但如果不懂得如何正确利用，可能反而会削弱个人的竞争力，适得其反，最终成为社会的蛀虫。\\
《七个习惯》的作者史蒂芬•柯维（Stephen
Covey）在书的最后一章里提到，我们可以传承给下一代并具有持久价值的，只有2样东西：根和翼。\\
``There are only two lasting bequests we can give our children - one is
roots, the other wings.''

\hypertarget{ux6839}{%
\subsection{根}\label{ux6839}}

回顾人类过去的130年间，虽然经历了两次世界大战浩劫，但也催生了许多前所未有的创新:

\begin{itemize}
\tightlist
\item
  电灯 （煤气灯）
\item
  汽车 （马车）
\item
  飞机
\item
  摩天大厦
\item
  半导体，集成电路
\item
  计算机
\item
  互联网
\item
  移动电话
\end{itemize}

我们现在使用的许多发明，实际上都是``发明者''基于前人的经验不断演化而来的成果。我在爱丁堡大学的古乐器博物馆见到过各种奇形怪状的古乐器，它们都是现代乐队中铜管和木管乐器的前身，如双簧管、单簧管、巴松管、圆号和小号等。就像达尔文先生在《物种起源》中描述的物种进化一样，经过优胜劣汰，最终演化成现代的乐器。\\
例如我们熟悉的口琴和手风琴，其实是欧洲人在中国传统乐器笙的基础上改造而来的。\\
人类创新并不仅限于前述领域，这只是我们看到的凤毛麟角，在医疗、电影、音乐、数学等众多领域，同样展现出了非凡的创造力。这些都是祖先留给我们的宝贵财富。\\
正因为每一代人传承并发展了前人的成果，我们才能取得现在的成就。这，便是我们的``根''。

\hypertarget{ux7ffc}{%
\subsection{翼}\label{ux7ffc}}

除了``根''的传承，祖先还留给了我们``翅膀''，赋予我们自由思考、打破负面基因传承的能力，而不仅仅是简单地将基因传递给下一代。\\
通过下面的两个故事，分别介绍两种不同的策略。

\hypertarget{ux98deux884cux5b9eux9a8c}{%
\subsubsection{飞行实验}\label{ux98deux884cux5b9eux9a8c}}

自古以来，人类就渴望飞翔，东西方对此都曾有过诸多尝试，但大多以失败告终。传说明代火箭专家陶成道在凳子上绑了47根火箭，并绑上风筝进行飞行试验，最终却因爆炸丧命。15世纪，意大利的达•芬奇也曾尝试模仿鸟类翅膀设计飞行器，但也仅仅是停留在图纸上，并未实现。

\begin{longtable}[]{@{}l@{}}
\toprule
\endhead
\vtop{\hbox{\strut 莱特兄弟（Wright
brothers）在1903年成功进行了机动飞机的首次飞行。在12月17日的第四次飞行中，他们在空中飞行了59秒，航程达到了259米。}\hbox{\strut \href{文件:404页替换图片.jpg}{300px}}\hbox{\strut 莱特兄弟并非工程学或数学方面的天才，但他们为了实现``飞翔''的梦想，不断从失败中吸取教训并持续改进。从1899年起，他们开始不断试验：}\hbox{\strut 1900年，他们首次试验滑翔机，但效果并不理想。}\hbox{\strut 1901年，他们将机翼面积从15平方米扩大到26平方米，并成功飞行了120米，但发现试验数据与计算数据不符。为了找出原因，他们制造了风洞来测量各种机翼设计的浮力，共试验了200种机翼设计。利用收集到的数据，他们开始制作第三架滑翔机。}\hbox{\strut 到1902年9月，莱特兄弟已完成了1000次试飞，其中空中飞行时间最长达到了26秒。正因为有了这些经验，他们设计了飞机的舵，使飞行员能够更好地操控飞机，随后他们开始设计使用四缸发动机和螺旋桨的机动飞机。}\hbox{\strut 1903年成功试飞后，他们继续完善设计。}\hbox{\strut 到1905年10月，飞机已经能够在空中飞行39分钟，并能在飞行员控制下完成转圈。}\hbox{\strut 飞行就是在不稳定的大气条件下操控一台同样不稳定的机器，充满了各种变数，因此极其困难。然而，莱特兄弟通过反复试验，发明了许多关键部件，如机翼、舵、发动机和螺旋桨等，并研究了它们之间是如何相互影响的。}}\tabularnewline
\bottomrule
\end{longtable}

莱特兄弟的飞行故事告诉我们，他们的成功绝非偶然。除了强烈的意志和拼搏精神，更重要的是善于从失败中汲取经验教训，并通过持续的试验和数据收集，不断优化，最终实现了突破。\\

\hypertarget{ux91cfux5b50ux6545ux4e8b-quantum-theory}{%
\subsubsection{量子故事 Quantum
Theory}\label{ux91cfux5b50ux6545ux4e8b-quantum-theory}}

1905年，年仅26岁的爱因斯坦发表了《狭义相对论》，这或许是大家耳熟能详的成果。然而，同年他还发表了一篇开创性的论文，解决了物理学中的另一大谜团：当金属片受到光照时会释放电子，但若光的频率低于某个临界值\(\nu_0\)，则无法释放电子。这一现象无法用传统物理学解释。\\
爱因斯坦基于物理学家马克斯·普朗克（Max
Planck）1900年提出的量子理论，提出了突破性的假设：光由光子（photon）组成，每个光子携带的能量\(E = h\nu\)
，其中\(h\)是普朗克常数。只要光子的能量超过金属原子的逸出能量
\(W\)（一个固定值），便可释放电子。这解释了为何光的频率必须高于\(\nu_0\)，而与光的强度无关。\\
在高中化学课中，我们可能好奇为何元素周期表中，氢（H）、锂（Li）、钠（Na）等位于表的左侧且具有类似的化学性质。波尔（Niels
Bohr）的原子模型为此提供了答案：电子分层排布，第一层最多容纳2个电子，第二层最多容纳8个电子。锂的电子排布为(2,1)，钠为(2,8,1)，它们的最外层均只有一个电子。这种电子层结构也解释了元素在三棱镜中折射出的特征光谱线：这些光谱线对应的是电子在不同能级之间跃迁时的能量变化。\\
然而，在更深入的有机化学学习中发现，原子的结构远比上述描述复杂。电子不再像行星围绕核心旋转，而是以电子云的形式分布在原子核周围，其位置无法精确确定，只能用概率描述，有些分布的形状很奇怪，如下图：

\url{文件:L01.jpg}

这一新的电子理论始于1925年，当时分为两派科学家分别提出改进模型：

\begin{enumerate}
\tightlist
\item
  海森堡Heisenberg（波尔的学生）和波恩Born提出用矩阵数学解释的方程式
\item
  薛定锷Schröedinger，基于德布罗意de
  Broglie提出电子可以用波动方程描述的理论，于1926年发表``薛定谔方程''（Schröedinger
  equation）\\
\end{enumerate}

虽然两派互不相让，觉得对方不够完善（其实两派理论之间没有冲突）。

\href{文件:科学家们2.jpg}{650px}

上图是1927在Brussel Savoy Conference的合照：\\
\#普朗克Max
PLANCK（前面第一排，左一）用了6年时间研究黑体热辐射，1900年提出量子理论（解决了以前光波理论无法解释的现象）

\begin{enumerate}
\tightlist
\item
  爱因斯坦Albert EINSTEIN（前面第一排，左四）基于量子论，提出光子
\item
  玻尔Niels
  BOHR（第二排，右一）提出原子架构模型，解释了元素周期表和氢的光谱
\item
  德布罗意Louis de BROGLIE（第二排，右三）电子有波的特性
\item
  海森堡Werner HEISENBERG（前面第一排，左四）与波恩Max
  BORN（第二排，右二）1925年提出用矩阵数学解释的方程式完善玻尔Bohr原子模型的不足
\item
  泡利Wolfgang PAULI（最后排，右四）也在1924年提出不相容原理Exclusion
  Principle,
  原子里的每层轨道最多只能有2个电子，如果2个，它们的自转必须相对\(m_s = + \frac{1}{2} or - \frac{1}{2}\)
  ，来解释氮(He Helium)的特征光谱线
\item
  薛定锷Erwin
  SCHROEDINGER（最后排，右六）1926年提出薛定谔方程,用波来解释原子架构
\item
  狄拉克P.A.M.
  DIRAC（第二排，右五）1928年结合相对论，和不相容原理，提出更完善的方程式
  Dirac Equation，整合了Heisenberg 的矩阵数学 和薛定谔方程
\end{enumerate}

从1900年到1928年，不到30年的时间内，主要由9位理论物理学家（后来都获得了诺贝尔奖）提出的量子理论
帮我们理清了原子结构，并提出不同的模型以解释元素和原子的特性。这些研究帮助人类更好地理解原子的物理特性，并能更准确地解释各种化学现象，如化学键的特性及为何各元素可归属为金属、绝缘体或半导体。\\
要了解4个量子数的应用，可以依据下表列出某元素的电子分布：

\[n =\] 能量级，例如：1，2，3

\[l=\] 轨迹类， 从0到n-1, 例如：0=s , 1=p , 2=d

\[m_l=\] 具体某轨迹 ，例如：\(l=2\)里 有5个不同轨迹
\(m_l = -2, -1, 0, +1, +2\)

例如，X(Cl Chlorine氯) 17个电子的分布如下：

\href{文件:!chlorine_Screenshot_2024-12-29_100317.2.jpg}{650px}

也可以使用周期表，从左到右，从上而下，逐一把电子加上，左边2栋是s层，最右边的6栋是p层，中间的10栋是d层。例如Chlorine，首先1层2个电子，得到He
；然后2s层2个电子，2p层6个电子，得到Ne
；然后3s层2个电子，最后3p层5个电子。

\href{文件:!periodic_table_Screenshot_2024-12-29_100607.jpg}{700px}

并可以利用Schrodinger公式，估计围绕原子核每层里电子分布的概率，得出以下s层，p层，和d层的三维分布。例如：p层的3种轨迹之间没有重叠；d层的5种也一样。

\href{文件:Quantum_numbers_Screenshot_2024-12-28_112328.jpg}{650px}

因为这些物理学家在20世纪初30年内所作出的贡献，大家更清楚每个元素的电子配置与周期表，这些已经是高中物理化学的基本教材。量子理论的重要性不亚于过去几百年间传统物理学家，如伽利略、牛顿到法拉第、麦克斯韦等的贡献。

\begin{description}
\tightlist
\item[]
= = = =
\end{description}

虽然我们都继承了祖先的根与翼，但并非每个人都能为人类的持续发展作出贡献。要实现这一目标，需要具备哪些关键元素呢？

\begin{longtable}[]{@{}l@{}}
\toprule
\endhead
\vtop{\hbox{\strut To be a powerful transition person, one must first
change his inner first.}\hbox{\strut 要成为有影响力的过渡者（a powerful
transition
person），无论采用哪类策略，我们必须首先进行由内而外的改变，只有这样的改变才能持久。}\hbox{\strut (只有在内部发生转变，任何人才能在自己可控的范围内进行创新和持续改善。）}\hbox{\strut 以上源自Covey先生的《七个习惯》。}}\tabularnewline
\bottomrule
\end{longtable}

\hypertarget{ux9b54ux672fux5e08ux4e54ux5e03ux65af1997ux5e74ux91cdux56deux82f9ux679c13ux5e74ux540eux518dux628aux82f9ux679cux5e26ux4e0aux9ad8ux5cf0}{%
\subsubsection{``魔术师''乔布斯1997年重回苹果，13年后再把苹果带上高峰}\label{ux9b54ux672fux5e08ux4e54ux5e03ux65af1997ux5e74ux91cdux56deux82f9ux679c13ux5e74ux540eux518dux628aux82f9ux679cux5e26ux4e0aux9ad8ux5cf0}}

1985年，乔布斯被迫离开苹果公司，随后创立了NEXT公司。直到1997年，苹果收购NEXT，乔布斯重返苹果（当时苹果正处于破产边缘）。到2010年，他成功将苹果公司带上美国市值最高的科技公司，但这期间也经历了多次失败。\\
乔布斯在他的传记中说：``是什么力量推动我前行？我认为大多数创新者（Creative
people）都希望向前辈致敬，感激他们对人类的贡献。\\
我并没有发明我用的语言或数学。我的食物并非由我自己生产，衣服也从未亲手制作过。我所做的一切都依赖于前人的积累。许多人都渴望回馈社会，在历史长河中留下属于自己的一笔。尽管我们无法创作鲍勃•迪伦（Bob
Dylan）的歌曲或汤姆•斯托帕德（Tom
Stoppard）的戏剧，但我们可以用自己的方式表达感激。试图用自己的才能表达内心感受，表达对前人的感激并做出贡献，这就是推动我的力量。''

\hypertarget{ux5171ux901aux70b9}{%
\subsection{共通点}\label{ux5171ux901aux70b9}}

为什么莱特兄弟、提出量子论的物理学家以及乔布斯，分别能够在航天、物理和电子产品领域实现突破性的重大贡献？他们的故事有两个共同点：

\begin{itemize}
\tightlist
\item
  专注------例如：

  \begin{itemize}
  \tightlist
  \item
    专心研究飞机
  \item
    解决物理难题，更好理解这个世界
  \item
    结合新技术研发针对个人市场的领先电子产品
  \end{itemize}
\item
  长期锲而不舍------从萌芽到成功的经历充满挑战，期间困难重重，必须保持内心的坚定。也正是这种坚定使他们成为具有影响力的过渡者，延续了人类发展的根与翼。以乔布斯为例，从被迫离开苹果到东山再起，历经25年，通过推出iMac、iPod、iPhone和iPad等创新产品，最终让苹果再创辉煌（这种公司级的突破，往往需要更长时间的积累与努力）。
\end{itemize}

从量子论的发展看到创新，除了要有超凡人才的专注和不断努力外，
还有以下特征：

\begin{itemize}
\tightlist
\item
  创新不是从零开始，而是基于之前的理论基础再突破
\item
  要抱着怀疑和批评的精神，看到不完善之处，例如：

  \begin{itemize}
  \tightlist
  \item
    爱因斯坦一直反对玻尔学派提出的量子力学不确定性原理，直言：``上帝不会掷骰子！''。他试图通过设计实验来证明这一理论的缺陷
  \item
    海森堡认为薛定谔的波函数方程（wave
    function）并不能完整描述原子中电子的所有特性，因而是不完善的
  \item
    薛定谔Schroedinger反而觉得海森堡Heisenberg的理论太抽象
  \end{itemize}
\end{itemize}

现在仍未形成一套完美的理论，能解释量子论的不足，但量子论已为我们带来了深远的影响，使我们能够理解许多过去无法解释的化学与物理现象。\\

\begin{description}
\tightlist
\item[]
= = =
\end{description}

``你的例子大多是超级明星级别的人物，``根与翼''这一概念未必能够完全适用于我们这些普通人。''\\
``为什么你认定自己不能成为有影响力的过渡者？''我问，然后讲我的故事：\\
`` 2000年，我被公司辞退，沦为失业汉，当时我39岁。\\
从2006年开始，香港政府鼓励软件开发公司使用CMMI模型做评估与改进，并提供资助。我便开始专注于软件开发CMMI过程改进的培训与评估，但香港市场很小，大陆市场巨大，但从未跑过，我从长三角地区开始，与大陆公司合作提供CMMI服务。2010年通过美国面试，成为CMMI高成熟度主任评估师。\\
2016年，我开始深入比较各种敏捷开发方法。2018年，正值美国PMI推出ACP个人证书，我开始与培训机构合作教授ACP课程。学习前辈在敏捷、精益和软件工程的贡献，成为我的软件工程的根，没有他们，我不可能从零自创出来。\\
不固步自封，因市场环境变化，从传统模式改变为精益思路，是我的翼，\\
这些都是我自身根和翼的经历，让我走过24年的独自创业之路，也帮助我出版第一本书。\\
所以任何人，只要把心态转过来，并能锲而不舍，都有机会。''\\
``好吧。但你说这些跟敏捷开发、持续改善有什么关系？我没有期望能成为爱因斯坦！''\\
``请问我们前面讨论过的
例如估算、需求调研、开发流程、持续集成、测试自动化、配置管理等。是否有改善的空间？''\\
``听你这样说，反观我们团队，我立马想到有很多地方需要改进，为什么你提议只针对减少缺陷引起的返工做改进？''\\
要理解背后的原因，先听听质量大师Dr Juran回顾过去质量改善的成功要素：

\begin{itemize}
\tightlist
\item
  All Improvement takes place Project by Project
  所有改进只能按项目逐个进行
\item
  Backlog of Improvement projects is huge 潜在的改进项目非常多
\item
  Quality Improvement does not come Free 质量改进并非免费
\item
  Major gains come from the Vital few projects
  主要效果（收益）来自少数重要项目
\item
  Quality improvement extends to All parameters
  质量改进可扩展到各类参数（如生产率、进度、安全性、环境等）
\end{itemize}

（以上源自Juran博士《Juran Quality Handbook》第五章质量改进。）\\
所以，如果从一开始同时对4到5个方向进行改进，反而会失去聚焦，难以尽早看到改进的效果。\\
但如果只针对降低缺陷引起的返工，目标明确，不仅让大家更容易接受，也更容易扩展到各个过程，例如需求、开放、产品集成，甚至配置管理。\\
因为能否引起高层对改进的兴趣非常重要，所以与企业领导首次见面时，我都会讲述团队应如何分析迭代缺陷数据，以此开始定量管理。\\

\hypertarget{tedux6f14ux8bb2}{%
\subsection{TED演讲}\label{tedux6f14ux8bb2}}

今年被质量经理邀请参加他们的质量沙龙，但只有30分钟，我就用了TED演讲思路，用了17分钟讲软件开发公司的常见质量问题和改进措施。

\begin{longtable}[]{@{}l@{}}
\toprule
\endhead
\begin{minipage}[t]{0.97\columnwidth}\raggedright
你觉得下面4种工作，哪一类工种占的工作量最多？

\begin{enumerate}
\tightlist
\item
  编码与代码设计。
\item
  交付后的所有工作，包括维护、更新与缺陷修正。
\item
  交付前的评审、静态扫描、测试与缺陷修正。
\item
  项目管理与监控。
\end{enumerate}

从2012年Capers
Jones对美国软件公司的统计中看到，编码只排第二（占开发工作量的25\%），而测试、评审和缺陷修改的总工作量则排第一（占30\%）。你选1.编码，其实大部分耗时的不是写代码，而是修正错误代码。软件开发有其特点，缺陷发现越晚，修复所需的工作量就越大。缺陷如果在早期被发现，可能只需要花费半小时或1小时进行修复。原因很简单，当已经集成为一个十几万行的系统，里面有很多组件。如果在完成集成后使用时才发现BUG，首先很难找到出问题的代码（原因很简单，当已经集成为一个十几万行的系统，里面有很多组件）；然后，为了验证代码修正是否有效，必须先测试单元模块是否通过，再测试集成后是否正常，最后还要通过系统测试才能交付给客户。相对而言，若能在自测、集成测试和评审阶段时尽早发现并解决问题，将节省大量时间。\\
然而，大部分缺陷都是在系统测试或集成测试时暴露的，有些产品线的缺陷数甚至达到四位数。如果能够将这些缺陷减少一半，将为公司节省近一半的工作量，同时也能改善产品质量，对吗？\\
道理大家都知道，但为什么没有改善？怎样才能改变现状，提前发现缺陷？\\
这种改进绝对不是通过下达一个命令或领导开会讲讲就能做到的。我先讲一个故事：在一家上万人的大型公司中，质量经理表示公司一直非常注重量化管理，为此制定了各种KPI考核团队，每月由主管分析数据。开始时确实有些改善，但过了一段时间后，效果就不明显了，甚至发现数据作假的情况。所以，我告诉质量经理，如果希望实现量化管理，就必须从团队提升开始，而不是仅靠总体数据分析和制定公司级改进策略。度量与分析最主要的作用在于帮助团队做好根因分析，制定纠正措施，并在下一轮迭代中进行提升。\\
例如，你们现在按2到4周进行一次迭代发布，确实也有一些团队在复盘，但他们不理解纠正措施与改正的区别，通常只是针对个别缺陷进行修复，结果仍会因为没有修复系统的基本问题而导致缺陷和事故重现。然而，如果整个团队在每轮迭代中分析数据、找出根因并制定纠正措施，并在下一个迭代中验证这些措施的有效性，就能养成良好的习惯。最重要的是，领导要给团队足够的时间进行复盘。要完成这种根因分析的复盘通常需要3个小时，因为要收集数据，并进行分析和讨论。\\
我们可以先从分析缺陷排除率入手，帮助团队养成复盘分析根因的习惯。缺陷有真实数据，让团队分析为什么需求引入的缺陷这么多？引导团队使用``五个为什么''方法找出根因。同时，团队成员需要了解当前的质量水平。通常，问需求人员自己的质量水平如何？他们会说需求评审中发现的BUG很少，就代表质量好，但这其实是一种误解。他们在评审过程中没有发现并不表示需求没有缺陷，等到客户发现时会更糟糕。\\
我们还可以利用一些预测模型和统计分析的方法。在迭代中，不仅要让团队制定纠正措施，还要设定定量化目标。例如需求评审应该发现多少缺陷？单元测试应该发现多少缺陷？如果未达到预期目标，说明可能还是会遗漏不少缺陷。通过设定这些目标，可以更好地推动团队进行前期的评审和单元测试。
如果团队开始养成分析根因并制定纠正措施的习惯，而领导又关注团队的持续改进，时不时地询问``下次迭代是什么时候？改进了多少？''，就可以扩大根因分析的范围。除了分析缺陷外，还可以把根因分析扩大到其他改进方向。\\
那么，如何开始呢？\\
只是把小孩扔进海里，他是学不会游泳的。因此，建议先进行培训，举办为期两天的工作坊，利用一些真实数据带领大家进行根因分析的互动练习，然后安排公司内部教练辅导团队，帮助他们在迭代中做好复盘。\\
\strut
\end{minipage}\tabularnewline
\bottomrule
\end{longtable}

\hypertarget{ux4e1cux5317ux864e}{%
\subsection{东北虎}\label{ux4e1cux5317ux864e}}

任何提升都需要内心思维的转变驱动。改变固有的思想和习惯并不容易。如果个人、团队或企业缺乏基于事实（数据）持续改善的思维，敏捷宣言、12条原则和各种方法（如SCRUM、XP、ACP），和相关的培训和技巧培训，最终都无法在敏捷开发中产生实际改进效果。（如首章A公司故事）\\
但如果个人和企业领导不满意当前的工作状态，内心相信必须改善质量、有危机意识，认为这工作重要，那么相关的方法与技巧反而次要。敏捷团队在改善过程中，或许不必单纯依赖敏捷原则和方法，甚至可以创造出更合适的方法与技巧。\\
团队和企业要健康成长依赖以下主要因素：\\
*高层要制订很高的量化目标，也要提供资源（包括给团队时间）。

\begin{itemize}
\tightlist
\item
  改进要有专注。

  \begin{itemize}
  \tightlist
  \item
    例如针对如何减少缺陷返工，尽早在评审和前期测试发现缺陷
  \end{itemize}
\item
  有能力与动力的团队成员。
\end{itemize}

有竞争才有进步，不知道自身的不足，缺乏危机感往往是改善的最大阻力。

\begin{longtable}[]{@{}l@{}}
\toprule
\endhead
\vtop{\hbox{\strut 东北虎属于濒危物种，国家组织人工饲养予以保护。}\hbox{\strut 专家曾试图将这些人工饲养的东北虎重新放归野外，但这一任务极为艰巨。它们是否能够在野外独立存活下来，仍然是一个未知数。有一部纪录片生动地讲述了动物保护组织在这一过程中所面临的各种挑战和困难。}\hbox{\strut 要确保老虎在野外能独立生存，它们必须能够自行觅食。在自然环境中，鹿是它们的主要猎物。然而，鹿极为敏捷，奔跑时速可达50公里，并具备很强的耐力。因此，专家需要评估老虎是否具备成功捕猎鹿的能力，否则无法将其放归野外。评估过程中，专家发现人工饲养的老虎因缺乏锻炼，其奔跑和狩猎能力明显不如野生老虎。为此，专家设计了多种锻炼方案，例如使用车辆牵引轮胎并悬挂鲜肉块以刺激它们奔跑。经过一系列锻炼后，老虎的体力有所恢复，但能否成功捕猎到鹿仍是未知数。接下来的训练是用汽车拖拽鹿的尸体吸引老虎奔跑。然而，有些老虎因为在平时的饲养过程中并没有吃过鹿肉，所以对此并不感兴趣。最终只有两三只老虎追着汽车跑，试图抓到鹿。最后，有两只年轻的老虎成功了。但将鹿的尸体作为奖励喂给获胜的老虎却并不容易------它们面对鹿的尸体无从下口，因为之前的食物都是由饲养员切好后进行喂食的。经过数小时尝试，这两只老虎才逐渐学会如何啃食鹿肉。}}\tabularnewline
\bottomrule
\end{longtable}

环境是影响团队动力的重要因素。看完这部纪录片，我不禁联想到一些本来充满抱负和动力的年轻程序员，他们同样深受环境影响。如果大家只顾加班赶进度，忽视质量，习惯于大量缺陷在测试阶段才暴露的生产方式，如何能期望这些年轻程序员注重代码质量呢？就像保护区试图将人工饲养的东北虎放归野外一样，这些新人在几个月后便适应了大环境，形成习惯，再改变就非常困难。\\
所以若想团队做好敏捷开发，为公司带来价值，首先管理层必须意识到当前的不足：如果不立马做改进，
便难以保持公司的竞争力。\\

\hypertarget{ux603bux7ed3}{%
\subsection{总结}\label{ux603bux7ed3}}

在软件开发领域，我们这一代人继承了自20世纪60年代以来硬件和软件的各种创新成果，为我们的发展提供了坚实的``根''。\\
敏捷宣言则指出了传统的瀑布式开发和计划驱动开发的局限性，为我们插上了``翼''。对于有能力的团队，这些方法提供了快速开发高质量产品的机会。\\

\begin{longtable}[]{@{}l@{}}
\toprule
\endhead
\vtop{\hbox{\strut ``请问你们这几年举办了多轮项目管理培训，有什么改进？''我与内部评估组成员一起按CMMI模型实践访谈交流时问某公司培训专员。}\hbox{\strut 她无言而对，我接着说：``虽然你们培训计划，考试打分，满意度调查，培训平台，内部讲师等都有，每月也定期汇报，但没有听到任何改善，高层只关注当年预算有没有用完，是否完成了领导要求的任务。千万不要误以为我们针对培训，每个岗位都应该把改进视为工作的一部分。''}}\tabularnewline
\bottomrule
\end{longtable}

所以无论\\
是个人、还是公司，\\
是改善、还是创新，\\
如果这个月没有比上个月有可衡量的提升，便应问自己为什么。

千万不要以为这个过程会一帆风顺，只要按计划执行就能达到目标。实际上，在实施过程中肯定会遇到各种无法预料的问题。只要明白当前的水平、分析根因、采取纠正措施并持续改进，就能不断接近目标，如丰田7个习惯中的：``持续改善永无止境''------就像龟兔赛跑中的乌龟一样。\\
丰田汽车在90年代初有9万员工，每年总共有450万条改进建议，平均每员工每年提出50条。但如果能坚持不断回顾、改善，应该能像乔布斯一样讲出你的动人故事。（详见附录C：与毕业生讲故事）

\hypertarget{ux53c2ux8003-references}{%
\section{参考 References}\label{ux53c2ux8003-references}}

\begin{enumerate}
\tightlist
\item
  Covey, Stephen R. \emph{The Seven Habits of Highly Effective People.}
  (Fireside Book , 1990)\\
\item
  Isaacson, Walter. \emph{Steve Jobs.} (Simon and Schuster Paperbacks,
  2011)\\
\item
  Juran, Joseph M. \emph{Juran's Quality Handbook} 5 ed.( McGraw - Hill,
  1999)\\
\item
  McEvoy, J. P. \emph{Introducing Quantum Theory} (Totem Books, 1996)\\
\end{enumerate}


